\documentclass[../main.tex]{subfiles}

\begin{document}
    Berdasarkan hasil praktikum pada Modul {\Roman{chapter}} tentang {\leftmark}, dapat diambil kesimpulan bahwa:
    \begin{enumerate}
        \item {
            {\Struct} merupakan tipe data dalam bahasa pemrograman C++. {\Struct} digunakan untuk mengelompokan tipe-tipe data agar lebih mudah dibaca dan digunakan. Selain {\struct}, \textit{class} juga dapat digunakan untuk membuat tipe data baru. Perbedaan {\struct} dengan \textit{class} ialah \textit{member} {\struct} bersifat \textit{public} secara \textit{default} sedangkan \textit{member} \textit{class} bersifat \textit{private} secara {\default}.
        }
        \item {
            Deklarasi {\struct} adalah proses mendefinisikan sebuah tipe data baru yang terdiri dari berbagai tipe data lain yang dikelompokkan bersama. {\Struct} digunakan untuk mengelompokkan data yang berhubungan ke dalam satu entitas yang lebih kompleks. Misalnya, sebuah {\struct} untuk merepresentasikan data mahasiswa bisa terdiri dari nama, NIM dan IPK.
        }
        \item {
            Akses {\struct} melibatkan operasi membaca atau menulis data di dalam {\struct} yang telah dideklarasikan dan diinisialisasi. Untuk mengakses anggota dari {\struct}, digunakan operator titik (\citecode{.}) untuk mengacu pada masing-masing elemen di dalam {\struct} tersebut. Misalnya, untuk mengakses nama dari {\struct} mahasiswa, digunakan sintaks \citecode{mahasiswa.nama}.
        }
        \item {
            {\Nested} {\struct} adalah penggunaan {\struct} di dalam {\struct} lain sebagai anggotanya. Hal ini memungkinkan pembuatan struktur data yang lebih kompleks dan terorganisir. Misalnya, sebuah {\struct} untuk representasi buku yang berisi {\struct} untuk pengarang dan penerbit, sehingga data terkait dapat dikelompokkan dengan lebih baik.
        }
        \item {
            {\sArray} \textit{of} {\struct} ialah sebuah {\sarray} yang setiap elemennya adalah {\struct}. Misalnya, {\sarray} untuk menyimpan data beberapa mahasiswa. {\Struct} \textit{of} {\sarray} ialah sebuah {\struct} yang memiliki anggota berupa {\sarray}. Misalnya, {\struct} yang memiliki {\sarray} untuk menyimpan nilai mata kuliah dan {\sarray} untuk menyimpan nama mata kuliah.
        }
    \end{enumerate}
\end{document}
